\section{Related-work matrix and extended discussion}
\label{app:related-matrix}

\begin{table*}[h]
\centering\scriptsize
\caption{\textbf{Positioning against the closest prior work.} ``Target'':
what is identified or bounded. ``Sharp'': bounds are attained by admissible
laws. ``Ranking'': a ranking metric (AUROC) is bounded. ``Learns'': a
predictor is trained against the identified set. ``$\Gamma$ testable'': the
sensitivity parameter can be bounded from data. ``Nuisance'': a coverage
statement under estimated nuisances is given.}
\label{tab:related}
\begin{tabular}{lllccccc}
\toprule
work & assumption & target & sharp & ranking & learns & $\Gamma$ testable & nuisance\\
\midrule
\citet{manski1990nonparametric} & none & mean outcome & yes & no & no & --- & no\\
\citet{rosenbaum2002observational} & $\Gamma$ on odds & test / effect & yes & no & no & no & no\\
\citet{lakkaraju2017selective} & leniency IV & failure rate at a target rate & point & no & no & --- & no\\
\citet{kallus2018confounding} & MSM($\Lambda$) & policy value & yes & no & policy & no & partial\\
\citet{dearteaga2018learning} & expert consistency & labels for censored units & no & no & yes & --- & no\\
\citet{coston2021characterizing} & MSM / IV & fairness metrics & yes & threshold & no & no & no\\
\citet{rambachan2022counterfactual} & IV / bounds & counterfactual risk & yes & no & no & partial & no\\
\citet{dorn2023sharp} & MSM($\Lambda$) & weighted means & yes & no & no & no & yes\\
\textbf{this paper} & $\dcsm(\Gamma)\equiv$ MSM($\Lambda^2$) & AUROC, risk, regret & yes & \textbf{yes} & \textbf{yes} & \textbf{lower bound} & \textbf{Thm.~5}\\
\bottomrule
\end{tabular}
\end{table*}

\subsection{Extended related work}
\label{app:related-full}

\Cref{tab:related} positions the paper. \emph{Selective labels.}
\citet{lakkaraju2017selective} named the problem and proposed contraction,
which uses leniency variation to estimate a model's failure rate at a target
acceptance rate; \citet{kleinberg2018human} applied it to bail;
\citet{dearteaga2018learning} exploit expert consistency. Contraction
identifies a point and needs the most lenient decision maker to accept at
least the target rate; we characterise everything the data cannot rule out and
use leniency to \emph{falsify} $\Gamma$ instead (\cref{prop:falsify}).
\emph{Sensitivity analysis.} $\dcsm(\Gamma)$ is a decision-censoring form of
Rosenbaum's bounds \citep{rosenbaum2002observational}, equivalent
(\cref{lem:msm}) to Tan's marginal sensitivity model \citep{tan2006distributional},
for which sharp bounds and inference exist for weighted averages and effects
\citep{zhao2019sensitivity, yadlowsky2018bounds, dorn2023sharp}; Manski's
bounds \citep{manski1990nonparametric} are the $\Gamma\to\infty$ endpoint.
\citet{kallus2018confounding} learn \emph{policies} robustly under the same
model, the closest antecedent to \cref{sec:learning}; their target is policy
value, ours predictive risk, and our inner supremum is closed-form without a
relaxation. \citet{coston2021characterizing} and
\citet{rambachan2022counterfactual} bound fairness metrics and counterfactual
risk under selective labels; none of this literature bounds a ranking metric
sharply, the gap \cref{thm:sharp-auc} fills, and none states per-unit coverage
under estimated nuisances, the gap \cref{thm:nuisance} addresses, using the
double-machine-learning product rate \citep{chernozhukov2018double}.
\emph{Missing data and DRO.} IPW and doubly robust estimation
\citep{robins1994estimation} require $\Gamma=1$; selection models
\citep{heckman1979sample} buy identification with joint normality. Formally
\eqref{eq:dcl} is a DRO problem \citep{duchi2021learning}, but the ambiguity
set is what the data fail to determine, not a modelling choice
(\cref{rem:dro}). \emph{Ranking and
uncertainty.} The rank identity is a population form of the Mann--Whitney
representation \citep{mann1947test, hanley1982meaning}, not previously used
to linearise AUROC in the outcome regression; the credal decomposition of
\cref{cor:uq} follows \citet{abellan2005difference} and contrasts with
ensembles \citep{lakshminarayanan2017simple, hullermeier2021aleatoric} that
assume the training labels sample the deployment law.
